\documentclass[10pt,a4paper]{amsart}
\usepackage[margin=19mm]{geometry}
\usepackage{amsmath,amssymb,mathtools}
\usepackage[scr=rsfs,cal=euler]{mathalfa}
\usepackage{xfrac}
\usepackage[unicode,hidelinks,bookmarks=false]{hyperref}
\newcommand{\ie}{i.e.\ }
\newcommand{\cf}{cf.\ }
\newcommand{\resp}{respectively\ }
\newcommand{\Subsection}{\subsection}
\providecommand{\nobreakdash}{\nobreak\hbox{-}}
\setlength{\emergencystretch}{2em}
\multlinegap0pt
\usepackage{cancel}
\usepackage{amssymb}
\usepackage{stackengine}
\usepackage{xcolor}
\newtheorem{proposition}{Proposition}
\title{Embedding Boolean ample monoids as full submonoids of Boolean inverse monoids}
\author{Mark V. Lawson}
\date{Source: arXiv:2604.07152v1 (CC BY 4.0)}
\begin{document}
\maketitle

\noindent\emph{Source and licence.} Embedding Boolean ample monoids as full submonoids of Boolean inverse monoids. Version arXiv:2604.07152v1 (CC BY 4.0), distributed by arXiv under the Creative Commons Attribution 4.0 International licence, http://creativecommons.org/licenses/by/4.0/, as recorded in the arXiv licence metadata for this exact version and shown on the abstract page https://arxiv.org/abs/2604.07152v1. Full licence notice: \texttt{licenses/arxiv-2604.07152-CC-BY-4.0.txt}.\par\smallskip

\noindent\textit{Editorial scope.} Counted unit: the article's proposition prop:homuz (source0.tex lines 705--710). The article's complete original proof of it is delivered with it (source0.tex lines 711--753, one \texttt{proof} environment of 43 lines). No mathematics is written, repaired or re-derived: the statement and the proof are the article's own lines, in the article's own notation, and no retained line is substituted or re-flowed.  No further source-local context is needed: every cross-reference of every retained line resolves inside the two delivered spans.  Nothing was omitted from any retained span, so the package carries no omission record.  No retained line contains a citation, so the wrapper carries no bibliography and no bibliography entry.  No retained line contains a figure environment, an image-inclusion command or drawing code, so no image is delivered and the package declares no asset dependency.  Editorial formatting only, in addition: an AMS \texttt{amsart} preamble that restates the article's own declarations and macros used by the retained lines, namely the environments \texttt{proposition} on separate counters, together with the standard packages those lines need.  The retained environments print in this document as Proposition 1 in order of appearance, whereas the article prints the same items with the numbering of its own full document.  The complete native source of the article as distributed in the arXiv e-print for this exact version is bundled unchanged as \texttt{sources/198070/source0.tex}.
\begin{proposition}\label{prop:homuz} Let $S$ be a full submonoid of a Boolean inverse monoid $T$.
We suppose also that $S$ is a Boolean ample monoid.
Then $\mathsf{C}(T)$ is a groupoid of fractions of $\mathsf{C}(S)$ iff
each non-zero element of $T$ can be written as a finite join of elements of the form $a^{-1}b$
where $a,b \in S$.
\end{proposition}

\begin{proof} We define first a functor $\iota \colon \mathsf{C}(S) \rightarrow \mathsf{C}(T)$
which is an embedding where $\iota (\mathsf{C}(S))$ is a wide subcategory of $\mathsf{C}(T)$.
This will depend only on the fact that $S$ is a full submonoid of $T$.
This will involve comparing prime filters in $S$ with prime filters in $T$.
Let $A$ be a prime filter of $S$.
Then 
$$A^{u} = \{t \in T \colon a \leq t \text{ for some } a \in A\}$$ 
is a prime filter in $T$ that contains elements of $S$.
Conversely, let $B$ be a prime filter in $T$ that contains elements of $S$.
Then 
$$B^{d} = B \cap S$$
is a prime filter in $S$.
Observe that $(A^{u})^{d} = A$ and $(B^{d})^{u} = B$.
We have therefore defined a bijection between the prime filters in $S$ and the prime filters in $T$ that contain elements of $S$.
We therefore define the functor $\iota$ by $A \mapsto A^{u}$.
We may now prove the proposition.

Suppose that each non-zero element of $T$ can be written as a finite join of elements of the form $a^{-1}b$ where $a,b \in S$.
Let $X$ be any prime filter of $T$.
Let $t \in X$ be any element, necessarily non-zero.
Then from the fact that we are dealing with a prime filter and given how $t$ can be written, there exists $a^{-1}b \in X$ for some $a,b \in S$.
Then $X = A \cdot B$ where $a^{-1} \in A$ and $b \in B$.
We may therefore  write $X = ((A)^{-1})^{-1} \cdot B$
where
$A = (a^{-1}\mathbf{r}(B))^{\uparrow}$
and
$B =(b \mathbf{d}(X))^{\uparrow}$.
Observe that $a \in A^{-1}$
and so both $A^{-1}$ and $B$ contain elements of $S$.
It follows that both $A^{-1}$ and $B$ are in the image of $\iota$.

To prove the converse,
let $t$ be any non-zero element of $T$.
By assumption, any prime filter $X$ of $T$ containing $t$ must have the form $A^{-1} \cdot B$ where $A$ and $B$ are 
prime filters of $T$ that have non-empty intersections with $S$.
Let $a \in A \cap S$ and let $b \in B \cap S$.
Thus $a^{-1}b \in A^{-1}  \cdot B = X$.
Now,  $X$ is a prime filter that contains both $t$ and $a^{-1}b$.
But any element less than $a^{-1}b$ has the form $a_{1}^{-1}b_{1}$ where $a_{1}, b_{1} \in S$;
this follows from the fact that $S$ is an order ideal of $T$.
Whence $X_{t} = \bigcup_{a^{-1}b \leq t, a,b \in S} X_{a^{-1}b}$.
By compactness $t = \bigvee_{i=1}^{n}a_{i}^{-1}b_{i}$, where $a_{i},b_{i} \in S$.
\end{proof}

\end{document}
